\begin{abstract}
Tropical peatlands in Southeast Asia represent globally critical terrestrial carbon repositories that have been destabilized by anthropogenic drainage and recurrent wildfire crises. In Central Kalimantan, Indonesia, the legacy of the 1990s Mega Rice Project (ex-PLG) excavated thousands of kilometers of artificial drainage canals across ombrotrophic peat domes. During drought episodes, desiccated peat matrices ignite, releasing catastrophic emissions of greenhouse gases and particulate matter. However, operational wildfire prediction is hindered by two fundamental dilemmas: the methodological pitfall of spatial autocorrelation in machine learning---where coordinate-dependent algorithms achieve ostensibly high accuracy via spatial memorization but collapse when deployed across novel regions---and polarized debates regarding corporate versus smallholder land tenure culpability. Here, we present a calibrated spatiotemporal investigation integrating 43,937 satellite-detected active fire observations from the NASA FIRMS constellation (VIIRS 375\,m and MODIS Aqua) during the October 2026 fire pulse, official peatland depth polygons (BBLSDLP/KHG 2019), projected canal drainage vectors in metric coordinates (UTM Zone 49S, EPSG:32749), ECMWF ERA5-Land atmospheric reanalysis, and cadastral spatial overlays encompassing 627 oil palm concessions (HGU) and 185 forestry concessions (PBPH). Using a rigorous two-model comparative design evaluated via 5-fold Spatial Block Cross-Validation (\texttt{GroupKFold} grouped by regency), we demonstrate that a baseline spatial memorization model (Model~A, incorporating geodetic coordinates) collapsed from a random holdout test ROC-AUC of 0.7956 to a near-random spatial CV ROC-AUC of 0.5399 ($F_1$-score: 0.0583, collapsing to 0.0000 in four of five folds). Conversely, our Pure Biophysical Model (Model~B, trained strictly without coordinates) retained robust predictive generalization across all five spatial blocks (holdout ROC-AUC: 0.7527; extreme fire recall: 65.82\%; regular fire precision: 92.1\%; spatial CV mean $F_1$-score: 0.1399, 2.4-fold higher than Model~A). Explainable artificial intelligence via TreeSHAP revealed that 7-day cumulative rainfall deficit ($<5.0$\,mm, mean absolute $\text{SHAP} = 0.0329$), Euclidean proximity to drainage canals ($<3.0$\,km, 0.0322), wind velocity (0.0254), consecutive dry spells ($>14$\,days, 0.0224), and deep peat domes ($>200$\,cm, 0.0209) are the primary determinants governing severe fire radiative power ($\text{FRP} \ge 12.0$\,MW). Cadastral overlay demonstrated that \textbf{75.4\%} (33,138 hotspots) occurred on unallocated Non-Concession land (\textit{Area Penggunaan Lain} / APL), while \textbf{16.6\%} (7,276 hotspots) and \textbf{8.0\%} (3,523 hotspots) fell within oil palm and forestry concessions, respectively. These findings resolve single-sector culpability narratives, mandate institutional governance over unallocated APL peatlands, and establish empirical thresholds (14-day dry spell trigger; $<3$\,km canal rewetting priority) for regional wildfire disaster mitigation.
\end{abstract}

\noindent\textbf{Keywords:} Tropical peatlands; Wildfire modeling; Spatial block cross-validation; Spatial memorization; Random Forest; TreeSHAP; Land tenure governance; Mega Rice Project; Fire Radiative Power; Central Kalimantan.
